
% =============================================================================
% Project : THE PHILOSOPHY OF VIABILITY
% Book    : Book I — Ontology of Viability
% Book ID : VS.B01
% Part    : Part I — Identity and Manifestation
% Part ID : VS.B01.P01
% Chapter : Chapter 08 — Comparative Philosophical Examination
% Chapter ID : VS.B01.P01.C08
%
% File :
%   2_MAINMATTER/PART_I_IDENTITY_AND_MANIFESTATION/
%   CHAPTER_08_COMPARATIVE_PHILOSOPHICAL_EXAMINATION/
%   sections/section_09_hume.tex
%
% Section :
%   Section 09 — David Hume
%
% Structural Identifier :
%   VS.B01.P01.C08.S09
%
% Purpose :
%   To reconstruct Hume's empiricist method, his account of
%   induction and causation, his analysis of identity and
%   personal identity, and the unresolved difficulty recorded
%   in the Appendix to the Treatise; then to examine the
%   ontological significance and limitations of these accounts
%   in relation to the commitments established in Part I.
%
% Scope :
%   The examination addresses impressions and ideas, the Copy
%   Principle, inductive inference, constant conjunction,
%   necessary connection, identity attribution, the succession
%   and association of perceptions, memory, the bundle
%   interpretation, the Appendix problem, and Hume's critique
%   of substance.
%
% Philosophical Restriction :
%   Hume's epistemological and psychological explanations are
%   not treated as ontological conclusions without an
%   independent argument. Identity is not equated with
%   invariance, persistence, continuity, memory, consciousness,
%   association, or the attribution of sameness.
%
% Claim Form and Epistemic Status :
%   Historical reconstruction: textually supported.
%   Interpretive claims: qualified where scholarly disagreement
%   remains.
%   Comparative objections and replies: argued assessments.
%   The ontological adequacy of Hume's account remains subject
%   to the limitations and open questions stated below.
%
% =============================================================================


\section{David Hume}
\label{sec:ch08_s09}

\index[name]{Hume, David}
\index[concept]{Identity}
\index[concept]{Personal Identity}
\index[concept]{Impressions}
\index[concept]{Ideas}
\index[concept]{Copy Principle}
\index[concept]{Causation}
\index[concept]{Constant Conjunction}
\index[concept]{Necessary Connection}
\index[concept]{Induction}
\index[concept]{Custom}
\index[concept]{Habit}
\index[concept]{Memory}
\index[concept]{Association}
\index[concept]{Bundle Theory}
\index[concept]{Substance}
\index[concept]{Imagination}
\index[concept]{Persistence}
\index[concept]{Continuity}


% =============================================================================
% 1. SCOPE AND COMPARATIVE ORIENTATION
% =============================================================================

\subsection*{Scope and Comparative Orientation}

% Structural ID: VS.B01.P01.C08.S09.SS01.P001
\phantomsection\label{para:VS.B01.P01.C08.S09.SS01.P001}

David Hume is an important interlocutor for the ontological
examination of Identity because his analysis places pressure
on the assumption that the identity attributed to an object
must correspond to an independently identifiable,
unchanging subject or principle underlying its successive
appearances. His account connects the attribution of identity
with relations among perceptions, the operation of memory,
and the transitions of the imagination. The resulting
challenge concerns not merely the identity of persons, but
also the distinction between a succession of related
appearances and the identity attributed to that succession.
\cite[1.4.6.5--14, SBN 253--258]{HumeTreatise}

% Structural ID: VS.B01.P01.C08.S09.SS01.P002
\phantomsection\label{para:VS.B01.P01.C08.S09.SS01.P002}

The examination must distinguish several questions that
Hume's philosophy brings into close relation without making
them identical. These include the origin of ideas, the
grounds of inference from experience, the origin of the idea
of necessary connection, the relations among successive
perceptions, and the attribution of identity to changing
objects or persons. A conclusion about one of these questions
does not automatically settle the others. In particular,
an explanation of how an idea of identity arises is not
yet an explanation of whether the being to which identity
is attributed possesses an ontological identity.

% Structural ID: VS.B01.P01.C08.S09.SS01.P003
\phantomsection\label{para:VS.B01.P01.C08.S09.SS01.P003}

The purpose of this section is therefore not to classify
Hume simply as a sceptic who denies the self, a naturalist
who accepts a constructed self, or an empiricist who reduces
reality to perceptions. Each of these descriptions captures
some aspects of an interpretive position, but none should
replace a reconstruction of the arguments and qualifications
in the Treatise. The section will distinguish Hume's
explicit claims, the interpretations advanced by scholars,
and the critical conclusions proposed here.

% Structural ID: VS.B01.P01.C08.S09.SS01.P004
\phantomsection\label{para:VS.B01.P01.C08.S09.SS01.P004}

The comparative question is whether Hume's analysis
successfully explains the identity of changing beings,
or whether it primarily explains why and under what
conditions human thought attributes identity to a succession
of related perceptions. This distinction is decisive for
Part I. The ontological status of Identity cannot be settled
merely by identifying the psychological conditions under
which an observer forms or sustains an idea of sameness.


% =============================================================================
% 2. HISTORICAL AND TEXTUAL LOCATION
% =============================================================================

\subsection{Historical and Textual Location}
\label{subsec:ch08-hume-historical}

% Structural ID: VS.B01.P01.C08.S09.SS02.P001
\phantomsection\label{para:VS.B01.P01.C08.S09.SS02.P001}

Hume's principal early philosophical work is
\textit{A Treatise of Human Nature}, published in stages
in 1739 and 1740. Book I, \textit{Of the Understanding},
contains the discussions most directly relevant to the
present examination: the origin of ideas, causal inference,
the idea of necessary connection, the external world,
and personal identity. The first
\textit{Enquiry concerning Human Understanding}, published
in 1748, revisits several central epistemological arguments
of the Treatise in a shorter and revised form.
\cite{HumeTreatise,HumeEnquiry,QuRadcliffeHume2026}

% Structural ID: VS.B01.P01.C08.S09.SS02.P002
\phantomsection\label{para:VS.B01.P01.C08.S09.SS02.P002}

The relation between the Treatise and the later
\textit{Enquiry} must not be assumed to be one of simple
verbatim repetition. Hume expressed dissatisfaction with
the execution of his early work and later directed readers
towards his subsequent philosophical writings. At the same
time, he also described the later work as a recasting and
simplification of principles that remained substantially
the same. The appropriate historical conclusion is that
the later text deserves independent examination, while
the precise degree of philosophical revision remains a
matter for interpretation.
\cite{QuRadcliffeHume2026}

% Structural ID: VS.B01.P01.C08.S09.SS02.P003
\phantomsection\label{para:VS.B01.P01.C08.S09.SS02.P003}

The question of personal identity requires particular
textual care. Hume's extended discussion appears in
Treatise I.iv.6. The first Enquiry does not reproduce
a corresponding chapter on personal identity. It would
therefore be inaccurate to present the Treatise's account
as though Hume had simply restated and reaffirmed it in
the Enquiry. The Appendix to the Treatise must also be
read alongside I.iv.6 because it records a difficulty
that Hume himself regarded as serious.
\cite[1.4.6, SBN 251--263; Appendix, SBN 633--636]{HumeTreatise}
\cite{QuRadcliffeHume2026}

% Structural ID: VS.B01.P01.C08.S09.SS02.P004
\phantomsection\label{para:VS.B01.P01.C08.S09.SS02.P004}

The present examination gives priority to the Treatise
where the argument concerns personal identity, and uses
the Enquiry where Hume reformulates his arguments about
induction and necessary connection. These works are not
treated as interchangeable textual witnesses. Their
differences, omissions, and revisions form part of the
historical evidence and must remain visible in the
interpretation.


% =============================================================================
% 3. EMPIRICISM, PERCEPTIONS, AND THE COPY PRINCIPLE
% =============================================================================

\subsection{Perceptions, Impressions, Ideas, and the Copy Principle}
\label{subsec:ch08-hume-copy-principle}

% Structural ID: VS.B01.P01.C08.S09.SS03.P001
\phantomsection\label{para:VS.B01.P01.C08.S09.SS03.P001}

Hume begins his account of the human mind by distinguishing
perceptions into two broad classes: impressions and ideas.
Impressions are characterized by greater force and
vivacity, while ideas are comparatively less forceful
perceptions, such as those involved in remembering,
thinking, or imagining. The distinction concerns the
manner in which perceptions are experienced; it does not
by itself establish that every perception has the same
function or that all forms of mental activity can be
reduced to one undifferentiated operation.
\cite[1.1.1.1--3, SBN 1--3]{HumeTreatise}

% Structural ID: VS.B01.P01.C08.S09.SS03.P002
\phantomsection\label{para:VS.B01.P01.C08.S09.SS03.P002}

The Copy Principle states, in its principal formulation,
that simple ideas are derived from and resemble simple
impressions. Hume employs this principle as a method
for clarifying ideas whose meaning appears obscure.
When a philosophical term is used, the investigator
should ask which impression supplies the original
content of the corresponding idea. If no appropriate
impression can be identified, the claim to possess a
clear and determinate idea becomes questionable.
\cite[1.1.1--2, SBN 1--10]{HumeTreatise}
\cite[2.5--9, SBN 17--22]{HumeEnquiry}

% Structural ID: VS.B01.P01.C08.S09.SS03.P003
\phantomsection\label{para:VS.B01.P01.C08.S09.SS03.P003}

The Copy Principle nevertheless requires qualification.
Hume himself discusses the possibility of a missing
shade of blue, a case in which a person may form an
idea of a shade that has not previously been experienced
as an impression. The example complicates any
unqualified claim that every simple idea must in fact
be copied from a corresponding prior impression.
The principle remains central to Hume's method, but
its exact scope and exceptions must not be concealed.
\cite[2.8, SBN 20--21]{HumeEnquiry}

% Structural ID: VS.B01.P01.C08.S09.SS03.P004
\phantomsection\label{para:VS.B01.P01.C08.S09.SS03.P004}

The ontological significance of this method must be
restricted. The Copy Principle concerns the origin,
content, and intelligibility of ideas. It does not
independently establish that whatever cannot be
represented through a corresponding impression does
not exist. The absence of an identifiable impression
may challenge the clarity of a proposed idea, but
a conclusion about the existence or non-existence
of its supposed object requires an additional argument.
The distinction between an account of cognition and
an account of Being must therefore remain intact.


% =============================================================================
% 4. INDUCTION AND THE LIMITS OF EXPERIENTIAL INFERENCE
% =============================================================================

\subsection{Induction, Custom, and the Limits of Reason}
\label{subsec:ch08-hume-induction}

% Structural ID: VS.B01.P01.C08.S09.SS04.P001
\phantomsection\label{para:VS.B01.P01.C08.S09.SS04.P001}

Hume's examination of induction concerns the grounds
on which we move from observed instances to unobserved
instances, and from past regularities to expectations
about future events. Such inferences are indispensable
to ordinary life and empirical inquiry, but their
justification cannot be established merely by
demonstrating a necessary connection between the
observed and the unobserved. The idea of an effect
does not follow deductively from the idea of its
supposed cause.
\cite[1.3.6.1--4, SBN 86--89]{HumeTreatise}

% Structural ID: VS.B01.P01.C08.S09.SS04.P002
\phantomsection\label{para:VS.B01.P01.C08.S09.SS04.P002}

Hume argues that the attempt to justify the uniformity
of nature through experience faces a circularity.
If the reliability of future regularities is justified
by appealing to the regularities observed in the past,
the argument already assumes that past patterns are
a reliable guide to future cases. If the same principle
is defended through demonstrative reasoning, it is
unclear how a matter of fact could be established
by a relation of ideas alone. The difficulty is
therefore not merely a lack of sufficient observations;
it concerns the form of the proposed justification.
\cite[4.2--4.23, SBN 26--35]{HumeEnquiry}

% Structural ID: VS.B01.P01.C08.S09.SS04.P003
\phantomsection\label{para:VS.B01.P01.C08.S09.SS04.P003}

Hume's positive explanation appeals to custom or habit.
Repeated experience of similar conjunctions produces
a disposition to expect one event upon the appearance
of another. The transition is not presented as a
demonstration that the future must reproduce the past.
It is an operation of human nature through which
expectation and belief arise. The account therefore
explains the psychological origin of inductive
expectation without claiming to provide a deductive
proof of the uniformity of nature.
\cite[5.1--5.22, SBN 39--55]{HumeEnquiry}

% Structural ID: VS.B01.P01.C08.S09.SS04.P004
\phantomsection\label{para:VS.B01.P01.C08.S09.SS04.P004}

The precise philosophical conclusion remains disputed.
Some interpretations emphasize the psychological
description of how the mind actually forms inductive
beliefs; others emphasize the limits of rational
justification. These readings are not identical.
The distinction matters because a psychological
explanation of belief does not automatically establish
that the belief is unjustified, while the failure
of a particular rational justification does not
by itself establish that the belief is false.
\cite[4.2--5.22]{HumeEnquiry}
\cite{Millican1995,QuRadcliffeHume2026}


% =============================================================================
% 5. CAUSATION, CONSTANT CONJUNCTION, AND NECESSARY CONNECTION
% =============================================================================

\subsection{Causation and the Idea of Necessary Connection}
\label{subsec:ch08-hume-causation}

% Structural ID: VS.B01.P01.C08.S09.SS05.P001
\phantomsection\label{para:VS.B01.P01.C08.S09.SS05.P001}

Hume's analysis of causation develops from the problem
of how we form the idea of a necessary connection
between events. In an individual instance, experience
presents one event followed by another. The observation
of succession does not, by itself, disclose an
additional necessary tie through which the second
event must follow from the first. Hume therefore
examines the origin of the idea of causal power
rather than assuming that its meaning is already
transparent.
\cite[1.3.14.1--15, SBN 155--163]{HumeTreatise}

% Structural ID: VS.B01.P01.C08.S09.SS05.P002
\phantomsection\label{para:VS.B01.P01.C08.S09.SS05.P002}

Repeated experience introduces the relation of
constant conjunction. Similar kinds of events are
observed to occur in recurring patterns of succession
and contiguity. The mind consequently acquires a
disposition to pass from the idea of one event to
the expectation of its usual attendant. The repeated
instances do not reveal a new sensible quality
within either event; rather, their repetition gives
rise to a new determination in the mind.
\cite[1.3.6.3--4, SBN 87--89; 1.3.14.16--23, SBN 163--166]{HumeTreatise}

% Structural ID: VS.B01.P01.C08.S09.SS05.P003
\phantomsection\label{para:VS.B01.P01.C08.S09.SS05.P003}

Hume identifies the impression from which the idea
of necessity is derived with the internal determination
or customary transition through which the mind moves
from one event to its usual attendant. The idea of
necessary connection is thus connected with the
experience of regular conjunction and the operation
of custom. In the Enquiry, Hume presents the same
central explanation through the customary transition
of the imagination.
\cite[1.3.14.20--31, SBN 164--170]{HumeTreatise}
\cite[7.26--30, SBN 73--78]{HumeEnquiry}

% Structural ID: VS.B01.P01.C08.S09.SS05.P004
\phantomsection\label{para:VS.B01.P01.C08.S09.SS05.P004}

The interpretation of this account is contested.
One family of readings emphasizes the priority
of constant conjunction and treats Hume as offering
a regularity-based account of causation. Other
readings argue that Hume's analysis does not simply
deny objective causal relations, and that his
discussion of the internal impression concerns
the origin of our idea of necessity rather than
the complete ontological status of causal relations.
The distinction between the two definitions of
cause in the Treatise and Enquiry is central to
this disagreement.
\cite[1.3.14.31, SBN 169--170]{HumeTreatise}
\cite[7.29, SBN 76--77]{HumeEnquiry}
\cite{StrawsonSecretConnexion1989,QuRadcliffeHume2026}

% Structural ID: VS.B01.P01.C08.S09.SS05.P005
\phantomsection\label{para:VS.B01.P01.C08.S09.SS05.P005}

It would therefore be inaccurate to summarize Hume's
position by saying without qualification that
causation exists only in the mind. Hume recognizes
that objects stand in relations of succession and
contiguity independently of the operations of
understanding. His difficulty concerns the claim
that an observable event discloses a necessary
productive power as a quality that can be directly
perceived or distinctly conceived. The relation
between the events, the inference made by the mind,
and the idea of necessity must be examined separately.
\cite[1.3.14.28--31, SBN 168--170]{HumeTreatise}

% Structural ID: VS.B01.P01.C08.S09.SS05.P006
\phantomsection\label{para:VS.B01.P01.C08.S09.SS05.P006}

The comparative importance of Hume's analysis lies
in the demand it places upon any ontology that
introduces causal, productive, or necessary relations.
It asks whether the proposed relation has been
distinguished from the observed regularities through
which it is identified, and whether the meaning of
necessity has been adequately explained. The present
section does not attempt to settle the general
ontology of causation, nor does it equate Hume's
account of causal inference with the account of
Expression developed elsewhere in Part I.


% =============================================================================
% 6. IDENTITY IN GENERAL: THE SAME OBJECT AND RELATED SUCCESSIONS
% =============================================================================

\subsection{Identity, Relation, and the Imagination}
\label{subsec:ch08-hume-identity-general}

% Structural ID: VS.B01.P01.C08.S09.SS06.P001
\phantomsection\label{para:VS.B01.P01.C08.S09.SS06.P001}

Hume's discussion of personal identity is preceded
by an analysis of identity and diversity in general.
He distinguishes the idea of an object that remains
invariable and uninterrupted through time from
the idea of a succession of distinct objects
connected by close relations. Considered precisely,
these are different ideas: a succession of related
objects is not numerically one uninterrupted object
merely because the objects are closely connected.
\cite[1.4.6.6--7, SBN 253--255]{HumeTreatise}

% Structural ID: VS.B01.P01.C08.S09.SS06.P002
\phantomsection\label{para:VS.B01.P01.C08.S09.SS06.P002}

Hume's central explanation concerns the tendency
of the imagination to confuse these two cases.
Relations among successive objects may make the
transition of thought from one to another so smooth
that the mind experiences the succession almost
as though it were contemplating one uninterrupted
object. The ease of this transition encourages
the attribution of identity to objects whose parts
or conditions have in fact changed.
\cite[1.4.6.6--7, SBN 253--255]{HumeTreatise}

% Structural ID: VS.B01.P01.C08.S09.SS06.P003
\phantomsection\label{para:VS.B01.P01.C08.S09.SS06.P003}

Hume describes the resulting tendency as a confusion
between identity and relation. When identity is
attributed to a succession of changing or interrupted
objects, the imagination may also feign an unknown
principle that supposedly connects the parts and
makes the succession appear uninterrupted. Hume
therefore questions whether the ordinary attribution
of identity in such cases is supported by an
invariable object or by a principle of union that
can be independently identified.
\cite[1.4.6.6--7, SBN 253--255]{HumeTreatise}

% Structural ID: VS.B01.P01.C08.S09.SS06.P004
\phantomsection\label{para:VS.B01.P01.C08.S09.SS06.P004}

The examples of a repaired ship, an organism, and
a river illustrate the range of Hume's analysis.
A ship may undergo extensive replacement of its
materials while continuing to be called the same
ship; an oak may undergo substantial material change
while its organized life continues; and a river
may remain the same river despite the continuous
replacement and movement of its waters. Hume uses
such examples to show how the relations among
successive parts and the nature of the object
affect the imagination's attribution of identity.
\cite[1.4.6.8--14, SBN 255--258]{HumeTreatise}

% Structural ID: VS.B01.P01.C08.S09.SS06.P005
\phantomsection\label{para:VS.B01.P01.C08.S09.SS06.P005}

The significance of these examples must not be
reduced to the claim that every changing object
is simply unreal or that every ordinary identity
judgment is false. Hume's analysis concerns the
distinction between strict identity and the
attribution of identity to successions whose parts
are related in particular ways. His account explains
why the imagination treats certain successions
as identical, but the ontological question of
whether the relevant being possesses an identity
cannot be settled merely by describing this
psychological operation.

% Structural ID: VS.B01.P01.C08.S09.SS06.P006
\phantomsection\label{para:VS.B01.P01.C08.S09.SS06.P006}

The difficulty for Part I is therefore precise.
If identity is treated as nothing more than the
smooth transition from one related condition to
another, the account risks reducing identity to
the relations through which sameness is attributed.
If, conversely, identity is treated as an additional
unchanging object behind the succession, it risks
reintroducing the kind of unexplained principle
that Hume criticizes. The ontological alternative
must avoid both reductions without merely repeating
that identity and relation are different concepts.
\SectionRef{sec:being-and-identity}
\SectionRef{sec:identity-and-determination}
\SectionRef{sec:identity-and-persistence}


% =============================================================================
% 7. PERSONAL IDENTITY, PERCEPTIONS, AND MEMORY
% =============================================================================

\subsection{Personal Identity and the Succession of Perceptions}
\label{subsec:ch08-hume-personal-identity}

% Structural ID: VS.B01.P01.C08.S09.SS07.P001
\phantomsection\label{para:VS.B01.P01.C08.S09.SS07.P001}

Hume begins his discussion of personal identity by
criticizing the claim that we possess an uninterrupted
and intimate consciousness of a simple, perfectly
identical self. He asks which impression could give
rise to the idea of such a self. Since particular
impressions and perceptions change and succeed one
another, none appears to provide the constant and
invariable impression that the proposed idea would
require.
\cite[1.4.6.1--3, SBN 251--252]{HumeTreatise}

% Structural ID: VS.B01.P01.C08.S09.SS07.P002
\phantomsection\label{para:VS.B01.P01.C08.S09.SS07.P002}

Hume then describes the mind through a succession
of perceptions. His theatrical analogy is often
quoted, but it must be read with its immediate
qualification: he warns that the comparison of
the theatre must not mislead us, because the
successive perceptions are all that he claims
to find in the mind, and he does not identify
a separate place or substance in which they
appear.\footnote{Hume's qualification is essential to the analogy: ``The comparison of the theatre must not mislead us.'' The passage does not establish that the mind is literally a theatre or that a further bearer has been discovered. \cite[1.4.6.4, SBN 252--253]{HumeTreatise}}

% Structural ID: VS.B01.P01.C08.S09.SS07.P003
\phantomsection\label{para:VS.B01.P01.C08.S09.SS07.P003}

An important qualification follows from Hume's
distinction between personal identity as it concerns
thought or imagination and personal identity as
it concerns the passions and the concern taken
in oneself. His principal analysis in this section
focuses on the former. It would therefore be
inaccurate to treat every dimension of personal
identity, self-concern, and practical involvement
as though Hume had reduced all of them to one
single operation of the imagination.
\cite[1.4.6.5, SBN 253]{HumeTreatise}

% Structural ID: VS.B01.P01.C08.S09.SS07.P004
\phantomsection\label{para:VS.B01.P01.C08.S09.SS07.P004}

Hume identifies resemblance, contiguity, and causation
as relations that connect successive perceptions
and facilitate the imagination's transition from
one to another. These relations do not make the
perceptions numerically identical. Rather, they
explain why a succession of distinguishable
perceptions may be contemplated as though it
constituted one continuing object. The attribution
of identity is thus connected with the relations
among perceptions and with the operation of
imagination.
\cite[1.4.6.6--7, SBN 253--255]{HumeTreatise}

% Structural ID: VS.B01.P01.C08.S09.SS07.P005
\phantomsection\label{para:VS.B01.P01.C08.S09.SS07.P005}

Memory has a distinctive role in this account.
It makes past perceptions available to reflection
and allows relations of resemblance among them
to become apparent. Hume also argues that memory
contributes to the production of the relation
of resemblance through which the imagination
moves among perceptions. It is therefore not
adequate to describe memory as either a passive
record that simply discovers a fully formed
personal identity or as the sole condition that
creates every aspect of identity.
\cite[1.4.6.16--18, SBN 259--261]{HumeTreatise}

% Structural ID: VS.B01.P01.C08.S09.SS07.P006
\phantomsection\label{para:VS.B01.P01.C08.S09.SS07.P006}

Causation extends the account beyond the perceptions
that are directly remembered. Hume argues that
causal relations connect perceptions into a
succession in which earlier and later perceptions
influence one another. Memory reveals part of
this chain, while causal reasoning permits the
imagination to extend the attributed identity
to periods and actions that are no longer
remembered. The account is consequently more
complex than the simple claim that a person
is identical with whatever the person currently
remembers.
\cite[1.4.6.18--20, SBN 260--262]{HumeTreatise}

\begin{center}

% =============================================================================
% Project : THE PHILOSOPHY OF VIABILITY
% Book    : Book I — Ontology of Viability
% Book ID : VS.B01
% Part    : Part I — Identity and Manifestation
% Chapter : Chapter 08 — Comparative Philosophical Examination
% Chapter ID : VS.B01.P01.C08
%
% File :
%   4_ASSETS/diagrams/
%   VS.B01.P01.C08.S09.SS07.P006.DG01.hume_attribution_of_personal_identity.tex
%
% Section / Chapter / Asset:
%   Diagram DG01 — Hume's Account of the Attribution of Personal Identity
%
% Structural Identifier:
%   VS.B01.P01.C08.S09.SS07.P006.DG01
%
% Asset Type : DG
%
% Purpose:
%   To represent the relations and psychological transitions
%   through which Hume explains the attribution of personal
%   identity to a succession of perceptions.
%
% Representational Limit:
%   The diagram represents Hume's explanatory account of
%   identity attribution. It does not establish that an
%   ontological identity is constituted by these psychological
%   relations or operations.
%
% Dependencies:
%   TikZ; positioning library.
%   Loaded by the project preamble.
%
% =============================================================================

\begin{figure}[htbp]
    \centering

    \begin{tikzpicture}[
        font=\small,
        node distance=0.72cm,
        process/.style={
            draw=blue!55!black,
            fill=blue!7,
            rounded corners=3pt,
            line width=0.7pt,
            text width=8.2cm,
            minimum height=0.85cm,
            align=center,
            inner sep=6pt
        },
        central/.style={
            process,
            fill=blue!14,
            font=\small\bfseries
        },
        qualification/.style={
            draw=black!45,
            fill=black!3,
            rounded corners=2pt,
            text width=11.2cm,
            align=center,
            font=\footnotesize,
            inner sep=7pt
        },
        flow/.style={
            ->,
            draw=black!65,
            line width=0.8pt
        }
    ]

    \node[process] (perceptions) {
        \textbf{Distinct and Successive Perceptions}\\
        Impressions and ideas occurring in succession
    };

    \node[process, below=of perceptions] (relations) {
        \textbf{Relations among Perceptions}\\
        Resemblance \quad Contiguity \quad Causation
    };

    \node[process, below=of relations] (memory) {
        \textbf{Memory and Custom}\\
        Recollection of past perceptions and the formation
        of habitual transitions
    };

    \node[process, below=of memory] (imagination) {
        \textbf{Easy Transition of the Imagination}\\
        Related perceptions are contemplated as though
        they belonged to one uninterrupted object
    };

    \node[central, below=of imagination] (attribution) {
        Attribution of Personal Identity
    };

    \draw[flow] (perceptions) -- (relations);
    \draw[flow] (relations) -- (memory);
    \draw[flow] (memory) -- (imagination);
    \draw[flow] (imagination) -- (attribution);

    \node[qualification, below=0.85cm of attribution] (limit) {
        \textbf{Interpretive qualification:}
        This sequence explains how identity is attributed
        in Hume's account. It does not, by itself, establish
        that the ontological identity of a person is identical
        with the relations or psychological operations shown.
    };

    \end{tikzpicture}

    \caption{Hume's account of the attribution of personal
    identity through relations among successive perceptions,
    memory, custom, and the imagination.}
    \label{fig:VS.B01.P01.C08.S09.SS07.P006.DG01}
\end{figure}

% =============================================================================
% END OF DIAGRAM DG01
% =============================================================================

\end{center}

% Structural ID: VS.B01.P01.C08.S09.SS07.P007
\phantomsection\label{para:VS.B01.P01.C08.S09.SS07.P007}

The distinction between the relations among perceptions
and the identity attributed to their succession
is crucial. Hume explains how memory, resemblance,
and causal relations contribute to the formation
and extension of the idea of personal identity.
Whether these relations constitute the ontological
identity of a person, or instead explain the
attribution of identity to a succession of
perceptions, remains a further question. The
psychological explanation cannot be treated
as an ontological conclusion without an
additional argument.


% =============================================================================
% 8. THE APPENDIX AND THE UNRESOLVED PRINCIPLE OF CONNECTION
% =============================================================================

\subsection{The Appendix: Hume's Unresolved Difficulty}
\label{subsec:ch08-hume-appendix}

% Structural ID: VS.B01.P01.C08.S09.SS08.P001
\phantomsection\label{para:VS.B01.P01.C08.S09.SS08.P001}

The Appendix to the Treatise is indispensable for
a responsible interpretation of Hume's account
of personal identity. After presenting his analysis
of perceptions and their relations, Hume returns
to the problem of explaining the principle that
connects distinct perceptions and makes us attribute
simplicity and identity to them. He acknowledges
that his account is defective and that he cannot
discover a satisfactory explanation of the
principles that unite successive perceptions
in thought or consciousness.
\cite[Appendix, paragraphs 19--21, SBN 635--636]{HumeTreatise}

% Structural ID: VS.B01.P01.C08.S09.SS08.P002
\phantomsection\label{para:VS.B01.P01.C08.S09.SS08.P002}

The difficulty is not merely that Hume has failed
to identify one more relation among perceptions.
His problem concerns the intelligibility of the
connection itself. If perceptions are distinct
existences, Hume argues, they form a whole only
through some connection. Yet the human understanding
does not discover a real connection among distinct
existences; it experiences a determination of
thought to pass from one perception to another.
The explanation of the unity attributed to the
succession therefore appears to require a principle
that the account cannot satisfactorily explain.

% Structural ID: VS.B01.P01.C08.S09.SS08.P003
\phantomsection\label{para:VS.B01.P01.C08.S09.SS08.P003}

The scholarly interpretation of this difficulty
remains unsettled. Some interpretations locate
the problem in a tension between the distinctness
of perceptions and the claim that they compose
one mind. Other interpretations emphasize the
nature of composition, the role of relations,
or the difficulty of explaining the unity of
consciousness without introducing a further
perception or principle. Garrett and Cottrell
offer influential but different analyses of
Hume's second thoughts. The disagreement must
remain visible rather than being resolved by
selecting one interpretation without argument.
\cite{Garrett1981,Cottrell2015,QuRadcliffeHume2026}

% Structural ID: VS.B01.P01.C08.S09.SS08.P004
\phantomsection\label{para:VS.B01.P01.C08.S09.SS08.P004}

The Appendix consequently prevents the bundle
interpretation from being presented as an
unqualified and completely settled solution.
Hume's main discussion advances an account of
the attribution of identity through related
perceptions, but his later qualification reveals
a serious difficulty concerning the principle
that unites them. The correct historical result
is not that Hume simply abandons every claim
made in I.iv.6, nor that the Appendix can be
ignored. It is that the proposed account remains
philosophically incomplete in a way Hume himself
recognized.


% =============================================================================
% 9. SUBSTANCE, SUBJECT, AND THE LIMITS OF REDUCTION
% =============================================================================

\subsection{Substance, Subject, and the Limits of Reduction}
\label{subsec:ch08-hume-substance}

% Structural ID: VS.B01.P01.C08.S09.SS09.P001
\phantomsection\label{para:VS.B01.P01.C08.S09.SS09.P001}

Hume's criticism of the idea of substance is connected
with his general demand that ideas be traced to
their impressions. He challenges the claim that
we possess a clear idea of a substance distinct
from the particular qualities or perceptions
associated with it. In the discussion of the
mind, this criticism is directed against the
assumption that a simple and enduring subject
can be identified independently of the perceptions
through which it is experienced.
\cite[1.4.5--6, SBN 232--263]{HumeTreatise}

% Structural ID: VS.B01.P01.C08.S09.SS09.P002
\phantomsection\label{para:VS.B01.P01.C08.S09.SS09.P002}

However, the absence of a clear idea of a substance
must not be converted without qualification into
the claim that no subject of experience exists.
The question whether Hume denies the existence
of subjects, denies only a particular conception
of an enduring subject, or suspends judgment
about what lies beyond perceptions is a disputed
interpretive issue. A theory of what can be
clearly conceived and an ontology of what exists
are related inquiries, but they are not identical.
\cite{QuRadcliffeHume2026,Cottrell2015}

% Structural ID: VS.B01.P01.C08.S09.SS09.P003
\phantomsection\label{para:VS.B01.P01.C08.S09.SS09.P003}

The present comparison must therefore avoid two
opposite simplifications. It must not attribute
to Hume a fully developed ontology of a substantial
self that his analysis explicitly challenges.
But it must also not present the bundle interpretation
as an uncontested proof that no subject or being
exists beyond the perceptions that can be identified
in reflection. The precise ontological reach
of Hume's argument remains a question requiring
further analysis.


% =============================================================================
% 10. THE STRONGEST DEFENCE OF HUME
% =============================================================================

\subsection{The Strongest Defence of Hume's Position}
\label{subsec:ch08-hume-defence}

% Structural ID: VS.B01.P01.C08.S09.SS10.P001
\phantomsection\label{para:VS.B01.P01.C08.S09.SS10.P001}

The strongest defence of Hume begins by emphasizing
the methodological discipline of his analysis.
He does not accept an obscure metaphysical term
merely because it has a long philosophical history.
He asks what idea is being expressed, how it is
formed, and what experience could make it intelligible.
This demand places a legitimate burden on any
ontology that introduces Identity as a fundamental
notion. The ontology must explain the meaning
and function of its claims rather than relying
on an unexplained label.

% Structural ID: VS.B01.P01.C08.S09.SS10.P002
\phantomsection\label{para:VS.B01.P01.C08.S09.SS10.P002}

The defender may also argue that Hume does not
need to posit an additional entity behind the
succession of perceptions. If the relations among
perceptions and the natural transitions of
imagination explain why identity is attributed,
then the introduction of a further unchanging
bearer may appear explanatorily unnecessary.
The account is intended to avoid multiplying
entities beyond what the analysis can justify.

% Structural ID: VS.B01.P01.C08.S09.SS10.P003
\phantomsection\label{para:VS.B01.P01.C08.S09.SS10.P003}

A further defence is that Hume's discussion need
not be interpreted as a universal theory of every
form of identity. His analysis can be understood
as an examination of the conditions under which
human beings attribute identity to objects and
persons whose parts or perceptions change.
The fact that this account does not provide a
complete ontology of Being would not by itself
show that it fails at its intended task.

% Structural ID: VS.B01.P01.C08.S09.SS10.P004
\phantomsection\label{para:VS.B01.P01.C08.S09.SS10.P004}

This defence is substantial. It prevents the
comparison from criticizing Hume merely because
his account does not answer every question
addressed by Part I. It also exposes a genuine
burden for the Philosophy of Viability: the
distinction between Identity and the relations
through which identity is recognized or attributed
must be made intelligible without reintroducing
an unexplained substrate or treating Identity
as an additional object attached to Being.


% =============================================================================
% 11. THE STRONGEST OBJECTION FROM PART I
% =============================================================================

\subsection{The Strongest Objection from the Ontology of Viability}
\label{subsec:ch08-hume-objection}

% Structural ID: VS.B01.P01.C08.S09.SS11.P001
\phantomsection\label{para:VS.B01.P01.C08.S09.SS11.P001}

The strongest objection from Part I is that Hume's
account risks moving from an explanation of the
formation of an idea to a conclusion about the
ontological status of that which the idea concerns.
The fact that the imagination attributes identity
through relations among successive perceptions
does not by itself establish that the being under
consideration possesses no identity apart from
this attribution. The psychological conditions
of an identity judgment and the ontological
condition of the being judged must be distinguished.

% Structural ID: VS.B01.P01.C08.S09.SS11.P002
\phantomsection\label{para:VS.B01.P01.C08.S09.SS11.P002}

The objection is not that Hume must accept a
separate and unchanging substance underlying
every succession. Part I does not require Identity
to be an additional object, hidden bearer, or
independent property attached to Being. Being
and Identity are analytically distinguishable
while remaining ontologically inseparable.
Consequently, rejecting a substantial substratum
does not, without further argument, settle the
status of Identity itself.
\SectionRef{sec:being-and-identity}

% Structural ID: VS.B01.P01.C08.S09.SS11.P003
\phantomsection\label{para:VS.B01.P01.C08.S09.SS11.P003}

A related difficulty concerns the reduction of
identity to relations among successive conditions.
Resemblance, contiguity, causal connection, and
the ease of imaginative transition may explain
why a succession is treated as one object. Yet
none of these relations is automatically identical
with the identity of the being whose conditions
are related. If identity is reduced to one of
these relations, the account must explain why
that relation, rather than another, constitutes
the relevant identity.
\SectionRef{sec:identity-and-determination}
\SectionRef{sec:identity-and-persistence}

% Structural ID: VS.B01.P01.C08.S09.SS11.P004
\phantomsection\label{para:VS.B01.P01.C08.S09.SS11.P004}

The Appendix makes the difficulty more acute.
If the perceptions are distinct and their connection
cannot be satisfactorily explained, the account
must clarify how they form the relevant unity
or whole. The difficulty is not resolved by
repeating that they are related, because the
nature and explanatory sufficiency of the relation
are precisely what require examination. Hume's
own dissatisfaction therefore reveals a question
that cannot be removed by presenting the bundle
account as a completed explanation.
\cite[Appendix, paragraphs 19--21, SBN 635--636]{HumeTreatise}

% Structural ID: VS.B01.P01.C08.S09.SS11.P005
\phantomsection\label{para:VS.B01.P01.C08.S09.SS11.P005}

The objection must nevertheless remain proportionate.
It does not establish that Hume's account of
identity attribution is false, that relations
among parts are irrelevant, or that every changing
being possesses one simple and unchanging identity.
It establishes a narrower burden: an account
of the psychological attribution of identity
does not by itself settle the ontological question
of identity, and an account that reduces identity
to relations must explain the status and sufficiency
of those relations.


% =============================================================================
% 12. REPLY AND QUALIFICATION
% =============================================================================

\subsection{Hume's Possible Reply and the Qualification of the Objection}
\label{subsec:ch08-hume-reply}

% Structural ID: VS.B01.P01.C08.S09.SS12.P001
\phantomsection\label{para:VS.B01.P01.C08.S09.SS12.P001}

Hume's strongest reply would be that the objection
assumes the existence of an ontological question
whose terms have not yet been made intelligible.
If Identity is distinguished from the relations
among perceptions, the critic must explain what
this distinction means and how it can be supported.
Otherwise, the appeal to an ontological identity
independent of every criterion of recognition
may appear to preserve the very obscure principle
that Hume's analysis seeks to eliminate.

% Structural ID: VS.B01.P01.C08.S09.SS12.P002
\phantomsection\label{para:VS.B01.P01.C08.S09.SS12.P002}

This reply cannot be dismissed by asserting that
Identity is fundamental. A foundational commitment
still requires conceptual discipline and an account
of its implications. The Philosophy of Viability
must show that Identity is not reducible to a
determination, a relation, an invariant feature,
a sequence of manifestations, or a psychological
act of attribution. It must also explain how
Identity can be ontologically inseparable from
Being without becoming an additional entity.
\SectionRef{sec:identity-and-determination}

% Structural ID: VS.B01.P01.C08.S09.SS12.P003
\phantomsection\label{para:VS.B01.P01.C08.S09.SS12.P003}

The reply from Part I is therefore limited but
substantive. The fact that the idea of Identity
is formed, recognized, or applied through relations
and cognitive operations does not establish that
Identity itself is constituted by those operations.
The origin of a concept, the conditions of its
application, and the ontological status of its
referent are distinct questions. Hume's analysis
provides a serious challenge to the intelligibility
of Identity, but it does not automatically answer
the ontological question merely by explaining
the formation of the idea.

% Structural ID: VS.B01.P01.C08.S09.SS12.P004
\phantomsection\label{para:VS.B01.P01.C08.S09.SS12.P004}

This reply does not yet provide a complete positive
solution to Hume's challenge. The present comparison
does not claim that Part I has already supplied
a final account of the connection between identity,
relations, and changing manifestations. It establishes
that these questions must be kept distinct and
that the burden of explanation remains. Where
the ontology requires a further argument, the
matter must remain open rather than being settled
by a verbal assertion.


% =============================================================================
% 13. BILATERAL EVALUATION
% =============================================================================

\subsection{Bilateral Evaluation}
\label{subsec:ch08-hume-evaluation}

% Structural ID: VS.B01.P01.C08.S09.SS13.P001
\phantomsection\label{para:VS.B01.P01.C08.S09.SS13.P001}

Hume's principal contribution to the comparative
problem is his analysis of the difference between
a succession of related conditions and the identity
attributed to that succession. He shows that
identity judgments can be shaped by resemblance,
contiguity, causation, memory, and the operations
of imagination. His analysis therefore prevents
a theory of identity from assuming that the
attribution of sameness is self-explanatory.

% Structural ID: VS.B01.P01.C08.S09.SS13.P002
\phantomsection\label{para:VS.B01.P01.C08.S09.SS13.P002}

The principal limitation is that the explanation
of identity attribution does not, without further
argument, settle the ontological status of Identity.
Hume's analysis raises the possibility that what
is ordinarily treated as one continuing object
may be a succession of related conditions. Yet
the account must still explain the status of
the relations and the connection through which
the succession is treated as a whole. The
Appendix makes clear that Hume himself did
not regard this difficulty as fully resolved.

% Structural ID: VS.B01.P01.C08.S09.SS13.P003
\phantomsection\label{para:VS.B01.P01.C08.S09.SS13.P003}

The relation between Locke and Hume is also
important. Locke distinguishes identity criteria
according to the kinds of things under consideration,
including the identity of masses, organisms,
human beings, and persons. Hume examines the
relations and imaginative operations through
which identity is attributed to changing objects
and persons. Both analyses challenge the
universalization of a single simple criterion,
but they differ in their accounts of consciousness,
memory, relation, and the status of the attributed
identity.
\cite[II.xxvii]{LockeEssay}
\cite[1.4.6]{HumeTreatise}

% Structural ID: VS.B01.P01.C08.S09.SS13.P004
\phantomsection\label{para:VS.B01.P01.C08.S09.SS13.P004}

The comparison therefore yields a limited
correspondence rather than equivalence. Hume
provides a powerful analysis of how identity
is attributed and how the idea of necessary
connection arises. The Philosophy of Viability
must nevertheless preserve the distinction
between the ontological identity of a being
and the relations or cognitive operations
through which that identity is determined,
recognized, or described. The comparison
supports further examination; it does not
establish the superiority or completeness
of either position.


% =============================================================================
% 14. CLAIM FORM, EPISTEMIC STATUS, AND OPEN QUESTIONS
% =============================================================================

\subsection{Claim Form and Epistemic Status}
\label{subsec:ch08-hume-status}

% Structural ID: VS.B01.P01.C08.S09.SS14.P001
\phantomsection\label{para:VS.B01.P01.C08.S09.SS14.P001}

The claims examined in this section do not all
possess the same logical form or justificatory
status. Historical statements about the contents
of the Treatise are textually supported claims;
interpretations of the relation between the
Treatise and the Enquiry are interpretive claims;
the objections developed from Part I are
arguments; and the unresolved status of the
Appendix is an open interpretive and philosophical
problem. These forms must not be collapsed into
one general label of acceptance or rejection.

% Structural ID: VS.B01.P01.C08.S09.SS14.P002
\phantomsection\label{para:VS.B01.P01.C08.S09.SS14.P002}

The functional or logical form of a claim must
be distinguished from its epistemic or justificatory
status. A historical reconstruction may be a
textually supported claim without becoming an
ontological axiom. A comparative thesis may be
argued without being conclusively established.
An open problem may be precisely formulated
without being treated as a proposition that
the book has already proved. This distinction
follows the methodological commitments established
for the chapter.
\SectionRef{sec:ch08-method-and-criteria}

% Structural ID: VS.B01.P01.C08.S09.SS14.P003
\phantomsection\label{para:VS.B01.P01.C08.S09.SS14.P003}

The following table summarizes the principal
claims and the limits of what they establish.
It is not a numerical evaluation of Hume's
philosophy. Its purpose is to prevent a claim
about the origin of an idea, the formation
of a belief, or the attribution of identity
from being silently converted into an ontological
conclusion.

\begin{center}

% =============================================================================
% Project : THE PHILOSOPHY OF VIABILITY
% Book    : Book I — Ontology of Viability
% Book ID : VS.B01
% Part    : Part I — Identity and Manifestation
% Chapter : Chapter 08 — Comparative Philosophical Examination
% Chapter ID : VS.B01.P01.C08
%
% File :
%   4_ASSETS/tables/
%   VS.B01.P01.C08.S09.SS14.P003.TL01.hume_claims_and_comparative_limits.tex
%
% Section / Chapter / Asset:
%   Table TL01 — Hume's Claims and Their Comparative Limits
%
% Structural Identifier:
%   VS.B01.P01.C08.S09.SS14.P003.TL01
%
% Asset Type : TL
%
% Purpose:
%   To distinguish Hume's principal arguments according to
%   their subject matter, explanatory function, and limits.
%
% Representational Limit:
%   The table is an analytical summary. It does not imply that
%   Hume's arguments form one completed or internally
%   uncontested ontological system.
%
% =============================================================================

\begin{table}[htbp]
\centering
\small

\caption{Hume's Principal Claims and Their Comparative Limits}
\label{tab:VS.B01.P01.C08.S09.SS14.P003.TL01}

\begin{tabularx}{\textwidth}{
    >{\raggedright\arraybackslash}p{2.25cm}
    >{\raggedright\arraybackslash}X
    >{\raggedright\arraybackslash}X
    >{\raggedright\arraybackslash}X
}
\toprule
\textbf{Problem}
&
\textbf{Hume's account}
&
\textbf{What it explains}
&
\textbf{What it does not establish}
\\
\midrule

Copy Principle
&
Simple ideas are derived from and resemble simple impressions,
subject to the qualifications raised by Hume himself.
&
The origin and intelligibility of ideas.
&
That whatever cannot be traced to an impression does not exist.
\\

Induction
&
Custom or habit explains the transition from observed
regularities to expectations about further instances.
&
The formation of inductive belief and expectation.
&
A deductive proof that the future must resemble the past,
or that every inductive belief is false.
\\

Causation
&
Constant conjunction and customary transition explain
the idea of necessary connection.
&
The formation of causal expectation and the idea of necessity.
&
An uncontested conclusion that causal relations exist
only in the mind or that objective causal relations
have been fully explained.
\\

Identity of objects
&
Relations among successive parts facilitate the imagination's
attribution of identity to changing or interrupted objects.
&
Why certain successions are ordinarily treated as one object.
&
That the ontological identity of an object is identical
with the psychological attribution of sameness.
\\

Personal identity
&
Resemblance, contiguity, causation, memory, and imagination
contribute to the attribution of identity to successive
perceptions.
&
The structure of Hume's account of personal identity.
&
A complete and uncontested explanation of the subject
or unity of experience.
\\

The Appendix
&
Hume acknowledges a difficulty in explaining the principle
that connects distinct perceptions.
&
The limits and unresolved tension within the earlier account.
&
That the bundle interpretation is a fully settled
and unproblematic final doctrine.
\\

\bottomrule
\end{tabularx}

\caption*{\footnotesize
The table distinguishes the explanatory task performed by
each argument from the broader conclusions that might be
drawn from it. The latter require additional argument and
must not be treated as automatic consequences of Hume's
textual claims.
}

\end{table}

% =============================================================================
% END OF TABLE TL01
% =============================================================================

\end{center}

% Structural ID: VS.B01.P01.C08.S09.SS14.P004
\phantomsection\label{para:VS.B01.P01.C08.S09.SS14.P004}

The principal open questions concern whether
Hume's account can explain the connection among
successive perceptions without presupposing
the unity it seeks to explain; whether the
relations that support identity attribution
are sufficient to explain the ontological
identity of changing beings; and whether
Identity can be defended without treating it
as either an additional substance or a product
of psychological attribution. These questions
remain open for further examination.


% =============================================================================
% 15. TRANSITION TO THE NEXT COMPARATIVE EXAMINATION
% =============================================================================

\subsection{Transition}
\label{subsec:ch08-hume-transition}

% Structural ID: VS.B01.P01.C08.S09.SS15.P001
\phantomsection\label{para:VS.B01.P01.C08.S09.SS15.P001}

Hume's examination has placed the relation between
identity, succession, connection, and attribution
under sustained pressure. The next comparative
examination will consider Kant, whose response
to the problems of experience, synthesis, and
the unity of consciousness introduces a different
account of the conditions under which perceptions
can be represented as belonging together. The
comparison must distinguish Kant's transcendental
conditions of experience from an ontological
account of Identity and must not assume that
the unity of apperception is equivalent to
personal identity.

% Structural ID: VS.B01.P01.C08.S09.SS15.P002
\phantomsection\label{para:VS.B01.P01.C08.S09.SS15.P002}

The result of the present section is therefore
a qualified one. Hume's analysis remains a
serious challenge to any account that treats
identity attribution, relation, and continuity
as self-explanatory. Yet his explanation of
the origin and operation of the idea of Identity
does not by itself settle the ontological status
of Identity. That question remains open to
the further comparative and philosophical
examination undertaken in this chapter.


% =============================================================================
% END OF SECTION 09 OF CHAPTER 08
% =============================================================================
